% !TEX root = task-6-4-1.tex
% Editable task bank. Append new BankTask blocks at the end of the document.
\documentclass[12pt,a4paper]{article}
\usepackage{physics-textbook}
\usepackage{graphicx,needspace}
\usetikzlibrary{decorations.pathmorphing}
\newif\ifShowSolutions
\ifdefined\TutorEdition
  \ShowSolutionstrue
\else
  \ShowSolutionsfalse % Change to \ShowSolutionstrue to display solutions.
\fi
\newcounter{banktask}
\newsavebox{\BankTaskBox}
\newcommand{\BankTaskContent}[3]{%
  \noindent\textbf{\thebanktask.}\quad #2\par
  \ifShowSolutions
    \Needspace{50mm}% Keep the solution heading with its opening text or diagram.
    #3\par
    \ifdefined\TutorEdition
      \ifcsname TutorSource#1\endcsname
        \nopagebreak\smallskip{\small\csname TutorSource#1\endcsname\par}%
      \fi
    \fi
  \fi
}
\newcommand{\BankTask}[3]{%
  \refstepcounter{banktask}\par\medskip
  \sbox{\BankTaskBox}{\begin{minipage}{\linewidth}%
    \BankTaskContent{#1}{#2}{#3}\end{minipage}}%
  \ifdim\dimexpr\ht\BankTaskBox+\dp\BankTaskBox\relax>\dimexpr\textheight-20mm\relax
    \Needspace{35mm}\BankTaskContent{#1}{#2}{#3}%
  \else
    \Needspace{\dimexpr\ht\BankTaskBox+\dp\BankTaskBox\relax}%
    \noindent\usebox{\BankTaskBox}\par
  \fi
  \medskip
}
\begin{document}
\BookSubsection{6.4.1}{Спектр испускания}


% Source sdamgia-10650: https://phys-ege.sdamgia.ru/problem?id=10650
\BankTask{sdamgia-10650}{
Электрон в атоме водорода переходит с энергетического уровня с номером n = $2$ на энергетический уровень с n = $1$. Чему равен модуль импульса испущенного при этом фотона? (В $10^{-29}\,\text{кг}\cdot\text{м/с}$; округлите до целого числа.)
}{
\smallskip\noindent\textbf{Решение}\par
Уровни энергии электрона в атоме водорода задаются формулой $E_n=-13{,}6\,\text{эВ}/n^2$, где $n=1,2,3,\ldots$ При переходе атома из состояния $E_{n}$ в состояние $E_{m}$ атом испускает фотон с энергией $E=E_n - E_m.$ При переходе со второго на первый энергетический уровень атом испустит фотон с энергией

\[E=E_2 - E_1=13{,}6\,\text{эВ}\cdot(1-\frac14) =10{,}2\,\text{эВ}.\]

Импульс фотона связан с его энергией как

\[p=\frac{E}{c}=\frac{10{,}2\cdot1{,}6\cdot10^{-19}\,\text{Дж}}{3\cdot10^8\,\text{м/с}}=544\cdot10^{-29}\,\text{кг}\cdot\text{м/с}.\]

Ответ: $544$.
}

% Source sdamgia-3811: https://phys-ege.sdamgia.ru/problem?id=3811
\BankTask{sdamgia-3811}{

\begin{center}\begin{tikzpicture}[book diagram,>=stealth]\draw[task axis,->](0,-.3)--(0,5.2)node[above]{$E$};\draw[thick](0,0)--(3.5,0);\node[left]at(0,0){$ E_0 $};\draw[thick](0,1.1)--(3.5,1.1);\node[left]at(0,1.1){$ E_1 $};\draw[thick](0,2.4)--(3.5,2.4);\node[left]at(0,2.4){$ E_2 $};\draw[thick](0,3.4)--(3.5,3.4);\node[left]at(0,3.4){$ E_3 $};\draw[thick](0,4.1)--(3.5,4.1);\node[left]at(0,4.1){$ E_4 $};\draw[->]( 0.7,0)--(0.7,3.4)node[midway,right]{$ 1 $};\draw[->]( 1.4,0)--(1.4,4.1)node[midway,right]{$ 2 $};\draw[->]( 2.1,2.4)--(2.1,0)node[midway,right]{$ 3 $};\draw[->]( 2.8,1.1)--(2.8,0)node[midway,right]{$ 4 $};\draw[dashed](0,5)--(3.5,5);\node[left]at(0,5){$0$};\end{tikzpicture}\end{center}
На рисунке изображена упрощённая диаграмма энергетических уровней атома. Нумерованными стрелками отмечены некоторые возможные переходы атома между этими уровнями. \TaskMatchingColumns{Процесс}{А) Поглощение света наибольшей длины волны\par Б) Испускание света наибольшей длины волны}{Стрелка}{1) 1\par 2) 2\par 3) 3\par 4) 4}\begin{center}\begin{tabular}{|c|c|}\hline А&Б\\\hline &\\\hline\end{tabular}\end{center}
}{
\smallskip\noindent\textbf{Решение}\par
Квант энергии поглощается при переходе на более высокий уровень, излучается при переходе на более низкий. При этом расстояние между уровнями обратно пропорционально длине волны: $\Delta E= \frac{hc}{\lambda} .$ Таким образом, поглощению света наибольшей длины волны соответствует энергетический переход $1$, а излучению света наибольшей длины волны отвечает переход $4$.

Ответ: $14$.
}

% Source sdamgia-4495: https://phys-ege.sdamgia.ru/problem?id=4495
\BankTask{sdamgia-4495}{

\begin{center}\begin{tikzpicture}[book diagram,>=stealth]\draw[task axis,->](0,-.3)--(0,5.2)node[above]{$E,\,\text{Дж}$};\draw[thick](0,4.6)--(3.5,4.6);\node[left]at(0,4.6){$ 0 $};\draw[thick](0,3.6)--(3.5,3.6);\node[left]at(0,3.6){$ -2\cdot10^{-19} $};\draw[thick](0,2.2)--(3.5,2.2);\node[left]at(0,2.2){$ -5\cdot10^{-19} $};\draw[thick](0,0.5)--(3.5,0.5);\node[left]at(0,0.5){$ -8\cdot10^{-19} $};\node[right]at(3.5,3.6){$E^{(3)}$};\node[right]at(3.5,2.2){$E^{(2)}$};\node[right]at(3.5,.5){$E^{(1)}$};\end{tikzpicture}\end{center}
На рисунке изображена схема низших энергетических уровней атома. В начальный момент времени атом находится в состоянии с энергией $E^{(2)}$. Согласно постулатам Бора с какой энергией данный атом может излучать фотоны? (В $10^{-19}\,\text{Дж}$.)
}{
\smallskip\noindent\textbf{Решение}\par
Согласно постулатам Бора, свет излучается при переходе атома на более низкие уровни энергии, при этом фотоны несут энергию, равную разности энергий начального и конечного состояний. Из приведенной здесь схемы видно, что фотон может излучиться только при переходе атома в состояние $E^{(1)} ,$ при этом его энергия будет равна $- 5 \cdot 10 ^{ - 19 } \,\text{Дж} - ( - 8 \cdot 10 ^{ - 19 } \,\text{Дж} ) =3 \cdot 10 ^{ - 19 } \,\text{Дж}.$

Ответ: $3$.
}

% Source sdamgia-3047: https://phys-ege.sdamgia.ru/problem?id=3047
\BankTask{sdamgia-3047}{
На рисунке изображены энергетические уровни атома и указаны длины волн фотонов, излучаемых и поглощаемых при переходах с одного уровня на другой. Какова длина волны фотонов, излучаемых при переходе с уровня $E_4$ на уровень $E_1$ если $\lambda_{13}=400\,\text{нм},$ $\lambda_{24}=500\,\text{нм},$ $\lambda_{32}=600\,\text{нм}$?
\begin{center}\begin{tikzpicture}[book diagram,>=stealth]\draw[task axis,->](0,-.3)--(0,5.2)node[above]{$E$};\draw[thick](0,0.3)--(3.5,0.3);\node[left]at(0,0.3){$ 1 $};\draw[thick](0,2.3)--(3.5,2.3);\node[left]at(0,2.3){$ 2 $};\draw[thick](0,3.7)--(3.5,3.7);\node[left]at(0,3.7){$ 3 $};\draw[thick](0,4.5)--(3.5,4.5);\node[left]at(0,4.5){$ 4 $};\draw[->]( 0.8,0.3)--(0.8,3.7)node[midway,right]{$ \lambda_{13} $};\draw[->]( 1.9,2.3)--(1.9,4.5)node[midway,right]{$ \lambda_{24} $};\draw[->]( 3.1,3.7)--(3.1,2.3)node[midway,right]{$ \lambda_{32} $};\end{tikzpicture}\end{center}
}{
\smallskip\noindent\textbf{Решение}\par
Частота фотона, испускаемого или поглощаемого атомом при переходе с одного уровня энергии на другой, пропорциональна разности энергий этих уровней. Поэтому имеем: $\nu_{41}=\nu_{43} + \nu_{31},$ $\nu_{43}=\nu_{42} - \nu_{32}.$
Отсюда (так как $\nu_{nk}=\nu_{kn}$): $\nu_{41}=\nu_{13} + \nu_{24} - \nu_{32}.$
Имеем:

$\nu_{13}= \frac{c}{\lambda_{13}} = \frac{3 \cdot 10 ^{ 8 }}{4 \cdot 10 ^{ - 7 }} =0{,}75 \cdot 10 ^{ 15 } \,\text{Гц}$
$\nu_{24}= \frac{c}{\lambda_{24}} = \frac{3 \cdot 10 ^{ 8 }}{5 \cdot 10 ^{ - 7 }} =0{,}6 \cdot 10 ^{ 15 } \,\text{Гц},$
$\nu_{32}= \frac{c}{\lambda_{32}} = \frac{3 \cdot 10 ^{ 8 }}{6 \cdot 10 ^{ - 7 }} =0{,}5 \cdot 10 ^{ 15 } \,\text{Гц},$
$\nu_{41}= ( 0{,}75 + 0{,}6 - 0{,}5 ) \cdot 10 ^{ 15 } \,\text{Гц}=0{,}85 \cdot 10 ^{ 15 } \,\text{Гц}.$
$\lambda_{41}= \frac{c}{\nu_{41}} = \frac{3 \cdot 10 ^{ 8 }}{0{,}85 \cdot 10 ^{ 15 }} \,\text{м}\approx 350\,\text{нм}.$

Ответ: $350\,\text{нм}$.
}

% Source sdamgia-4758: https://phys-ege.sdamgia.ru/problem?id=4758
\BankTask{sdamgia-4758}{
Уровни энергии электрона в атоме водорода задаются формулой $E_n= - \frac{13{,}6}{n ^2} \,\text{эВ},$ где $n = 1, 2, 3, ...$ . При переходе атома из состояния $E_2$ в состояние $E_1$ атом испускает фотон. Попав на поверхность фотокатода, фотон выбивает фотоэлектрон. Длина волны света, соответствующая красной границе фотоэффекта для материала поверхности фотокатода, 	$\lambda_{\text{кр}} = 300\,\text{нм}.$ Чему равна максимально возможная кинетическая энергия фотоэлектрона?
}{
\smallskip\noindent\textbf{Решение}\par
Согласно постулатам Бора, свет излучается при переходе атома на более низкие уровни энергии, при этом фотоны несут энергию, равную разности энергий начального и конечного состояний. Таким образом, испущенный фотон имел энергию

\[h\nu=E_2 - E_1= - 13{,}6 \,\text{эВ} ( \frac{1}{2 ^2} - \frac{1}{1 ^2} ) = 10{,}2\,\text{эВ}.\]

Согласно уравнению фотоэффекта, максимальная кинетическая энергия вылетающих фотоэлектронов связана с энергией фотона и работой выхода соотношением

\[h\nu=A + E_{\,\text{к}}.\]

Работа выхода связана с длиной волны красной границы соотношением: 

\[A= \frac{hc}{\lambda_{\,\text{кр}}}\]

Таким образом, максимально возможная кинетическая энергия фотоэлектрон равна

\[E_{\,\text{к}}=h\nu - \frac{hc}{\lambda_{\,\text{кр}}} = 10{,}2 \,\text{эВ} - \frac{4{,}1 \cdot 10 ^{ - 15 } \,\text{эВ} \cdot \,\text{с} \cdot 3 \cdot 10 ^{8} \,\text{м/с}}{3 \cdot 10 ^{ - 7 } \,\text{м}} = 6{,}1 \,\text{эВ}.\]

Ответ: $E_{\,\text{к}} = 6{,}1\,\text{эВ}$
}

% Source sdamgia-5388: https://phys-ege.sdamgia.ru/problem?id=5388
\BankTask{sdamgia-5388}{
Электрон, имеющий импульс $p=2\cdot10^{-24}\,\text{кг}\cdot\text{м/с}$ сталкивается с покоящимся протоном, образуя атом водорода в состоянии с энергией $E_{n}$ ($n = 2$). В процессе образования атома излучается фотон. Найдите частоту $\nu$ этого фотона, пренебрегая кинетической энергией атома. Уровни энергии электрона в атоме водорода задаются формулой $E_n= - \frac{13{,}6}{n ^2} \,\text{эВ},$ где $n=1,2,3,... .$
}{
\smallskip\noindent\textbf{Решение}\par
Поскольку кинетической энергией атома пренебрегаем, то вся кинетическая энергия электрона $E_0= \frac{p ^2}{2m_e}$ перейдет в энергию атома $E_2= - \frac{13{,}6}{2 ^2} \,\text{эВ}$ и энергию фотона.

Таким образом, по закону сохранения энергии энергия фотона равна 

\[E_{\,\text{фот}} =E_0 - E_2= \frac{p ^2}{2m_e} + \frac{13{,}6\,\text{эВ}}{2 ^2} =h\nu.\]

Отсюда находим частоту фотона:

\[\nu= \frac{p ^2 / ( 2m_e ) + (13{,}6\,\text{эВ})/2 ^2}{h}\]

 Тогда 

$\nu= \frac{4 \cdot 10 ^{ - 48 } \,\text{кг} \cdot \,\text{Дж}/ ( 2 \cdot 9{,}1 \cdot 10 ^{ - 31 } \,\text{кг} ) + 13{,}6/4 \cdot 1{,}6 \cdot 10 ^{ - 19 } \,\text{Дж/эВ} \cdot \,\text{эВ}}{6{,}62 \cdot 10 ^{ - 34 } \,\text{Дж} \cdot \,\text{с}} =$

$=4{,}14 \cdot 10 ^{ 15 } \,\text{Гц}\approx 4{,}1 \cdot 10 ^{ 15 } \,\text{Гц}.$

Ответ: $\nu=4{,}1 \cdot 10 ^{ 15 } \,\text{Гц}.$
}

% Source sdamgia-6513: https://phys-ege.sdamgia.ru/problem?id=6513
\BankTask{sdamgia-6513}{
Значения энергии электрона в атоме водорода задаются формулой $E_n= \frac{- 13{,}6\,\text{эВ}}{n ^2} ,n=1,2,3,....$ При переходе с верхнего уровня энергии на нижний атом излучает фотон. Переходы с верхних уровней на уровень с $n=1$ образуют серию Лаймана; на уровень с $n=2$  — серию Бальмера; на уровень с $n = 3$  — серию Пашена и так далее. Найдите отношение $\beta$ минимальной частоты фотона в серии Бальмера к максимальной частоте фотона в серии Пашена.
}{
\smallskip\noindent\textbf{Решение}\par
Частота фотона связана с его энергией равенством $h\nu=E,$ где $h$  — постоянная Планка. В серии Бальмера энергия фотона равна $E_n - E_2,$ где $n=3, 4, ....$ Аналогично в серии Пашена энергия фотона равна $E_n - E_3,$ где $n=4, 5, ....$ Минимальной частота фотона в серии Бальмера будет при условии перехода с $3$-го уровня, максимальной частота фотона в серии Пашена будет при переходе с самого высокого ($n = \infty$) уровня. Поэтому 

\[\beta = \frac{E_3 - E_2}{E_{\infty} - E_3} = \frac{\frac{1}{2 ^2} - \frac{1}{3 ^2}}{\frac{1}{3 ^2} - 0} =1{,}25.\]

Ответ: $1{,}25$.
}

\end{document}
